\section{The uncapped economy with unit elasticity}
\label{app:rewrite-uncapped-unit}

This appendix derives the rates and allocation ratios supporting
Proposition~\ref{prop:rewrite-uncapped-unit-bgp} in
Section~\ref{subsec:rewrite-uncapped-unit-bgp} and defines the normalized
variables used to prove local convergence. I use the uncapped RSI technology
\eqref{eq:rewrite-uncapped-unit-law} and the production function
\eqref{eq:rewrite-uncapped-unit-elasticities}.
In this appendix and its proofs, a superscript $*$ denotes the reference
BGP characterized in Section~\ref{subsec:rewrite-uncapped-unit-bgp},
distinguishing it from nearby equilibrium trajectories.

\subsection{An explicit balanced-growth equilibrium}

At unit elasticity, $s_X=\omega_X$ and $e_X=1-(1-\alpha)\omega_X$. The AI-demand and monopoly conditions,
equations~\eqref{eq:rewrite-ai-price} and~\eqref{eq:rewrite-monopoly-foc}, give
\begin{equation}
 \frac{p_XX}{Y}=(1-\alpha)\omega_X,\qquad
 \frac UY=(1-\alpha)^2\omega_X^2.
 \label{eq:rewrite-uncapped-unit-static-shares}
\end{equation}
Substituting $X=BU=(1-\alpha)^2\omega_X^2BY$ into production yields
\begin{equation}
 Y^{1-(1-\alpha)\omega_X}
 =K^\alpha(AL)^{(1-\alpha)\omega_L}
 \bigl[(1-\alpha)^2\omega_X^2B\bigr]^{(1-\alpha)\omega_X}.
 \label{eq:rewrite-uncapped-unit-production}
\end{equation}
On an exact balanced-growth path with positive constant $K/Y$, $C/Y$, and $M/Y$,
log differentiation of \eqref{eq:rewrite-uncapped-unit-production}
and~\eqref{eq:rewrite-uncapped-unit-law} gives the two growth restrictions
\eqref{eq:rewrite-uncapped-unit-production-growth}
and~\eqref{eq:rewrite-uncapped-unit-research-growth}.
Define
\begin{equation}
 \Delta\equiv(1-\alpha)(1-\eta-\omega_X).
 \label{eq:rewrite-uncapped-unit-feedback}
\end{equation}
For $\Delta>0$, their unique solution has
\begin{equation}
 g_Y^*=\frac{(1-\eta)(1-\omega_X)}{1-\eta-\omega_X}(n+\gamma),
 \label{eq:rewrite-uncapped-unit-output-growth}
\end{equation}
and $g_B^*$ as in \eqref{eq:rewrite-uncapped-unit-efficiency-growth}.
The equality $g_K=g_C=g_Y^*$ follows from constant capital and consumption ratios;
equations~\eqref{eq:rewrite-wage} and~\eqref{eq:rewrite-euler} then give
\eqref{eq:rewrite-uncapped-unit-wage-growth}
and~\eqref{eq:rewrite-uncapped-unit-return}.

Define the candidate ratios $k^*=K/Y$, $u^*=U/Y$,
$i^*=(\dot K+\delta K)/Y$, $m^*=M/Y$, and $c^*=C/Y$ by
\begin{align}
 k^*&=\frac{\alpha}{r^*+\delta},\qquad u^*=(1-\alpha)^2\omega_X^2,
 \label{eq:rewrite-uncapped-unit-capital-inference}\\
 i^*&=(g_Y^*+\delta)k^*,
 \label{eq:rewrite-uncapped-unit-investment}\\
 m^*&=\frac{u^*\eta g_B^*}{\rho-n+\eta g_Y^*},
 \label{eq:rewrite-uncapped-unit-research-share}\\
 c^*&=1-u^*-i^*-m^*.
 \label{eq:rewrite-uncapped-unit-consumption}
\end{align}
In this appendix, $k^*$ is a capital--output ratio, not capital per unit of
effective labor, and $i^*$ is gross physical investment as a share of output.
The research share comes from the developer's first-order
and costate conditions; the consumption share comes from market clearing.
These ratios are endogenous, not fixed saving or research rules.
The proof in Appendix~\ref{app:rewrite-proofs} verifies their positivity
and constructs the reference initial levels from the production and
research equations, for given initial population and labor-augmenting technology.
Along the BGP, $B_t=B_0^*e^{g_B^*t}\to\infty$,
$g_K=g_C=g_U=g_M=g_Y^*$, $g_X=g_Y^*+g_B^*$, and $g_q=g_Y^*-g_B^*$.

Since $n+\gamma>0$, equations
\eqref{eq:rewrite-uncapped-unit-production-growth}
and~\eqref{eq:rewrite-uncapped-unit-research-growth} require $\Delta>0$ for
this positive finite-rate BGP. If $\Delta=0$, they are inconsistent; if
$\Delta<0$, they imply $g_B<0$, contrary to the RSI technology. Thus
$\eta+\omega_X\geq1$ excludes this BGP, not every possible equilibrium
trajectory. The same restriction $\eta+\omega_X<1$ also suffices for
global developer optimality: the proof uses the invertible state transformation
$S=B^{(1-\eta)/\eta}$, in which the RSI technology is jointly concave in
$S,M$ and operating profit is concave in $S$. Finally, $\rho>n$ makes
both transversality expressions decline and the objective integrals converge.

\subsection{Equilibrium trajectories near the balanced-growth path}

The BGP supplies a reference for proving existence from nearby initial
stocks that need not be on that path. For the given exogenous paths of
population and labor-augmenting technology, compare each variable with its
value on the corresponding reference BGP at the same date. Define log deviations
\begin{equation}
 \xi_K=\log\frac K{K^*},\quad
 \xi_B=\log\frac B{B^*},\quad
 \xi_C=\log\frac C{C^*},\quad
 \xi_q=\log\frac q{q^*}.
 \label{eq:rewrite-uncapped-unit-deviations}
\end{equation}
Time subscripts are suppressed: the reference quantities in the denominators
follow the BGP rather than remain fixed at their initial values.
Substitution into the four equilibrium dynamics gives a smooth autonomous
system $\dot\xi=f(\xi)$ with $f(0)=0$, where
$\xi=(\xi_K,\xi_B,\xi_C,\xi_q)'$. Let $J^*=Df(0)$ denote its Jacobian.
The proof in Appendix~\ref{app:rewrite-proofs} derives its entries directly from
equations~\eqref{eq:rewrite-eq-kdot}--\eqref{eq:rewrite-eq-qdot}, with
$\psi=1$ and $\psi'=0$.
The normalized Jacobian is independent of $\chi$, $A_0$, and $N_0$;
changing research productivity changes the reference levels, not the
eigenvalues governing local convergence.

The parameter restrictions in Proposition~\ref{prop:rewrite-uncapped-unit-bgp}
imply two local properties. First, the linearized dynamics have exactly two stable
directions and no eigenvalue with zero real part. Writing $E^s$ for the real
invariant subspace associated with eigenvalues having negative real parts,
the first property is
\begin{equation}
 \operatorname{spec}(J^*)\cap i\mathbb R=\varnothing,
 \qquad \dim E^s=2.
 \label{eq:rewrite-uncapped-unit-local-spectrum}
\end{equation}
Second, these stable directions allow independent changes in the two
predetermined stocks. Let $V_S\in\mathbb R^{4\times2}$ be any real basis
matrix for $E^s$, and let $V_{KB}$ consist of its first two rows. Then
\begin{equation}
 \det V_{KB}\ne0.
 \label{eq:rewrite-uncapped-unit-local-projection}
\end{equation}
Nonsingularity does not depend on the choice of basis. It allows the stable
manifold to select $C_0,q_0$ for every sufficiently close pair $K_0,B_0$,
rather than only for stocks on a curve. Both properties follow analytically
from the proposition's parameter restrictions; neither is an additional
assumption or a calibration-specific test. The complete argument, including both transversality
conditions and global agent optimality, is in Appendix~\ref{app:rewrite-proofs},
\hyperref[proof:rewrite-uncapped-unit-local]{Proof of
Proposition~\ref*{prop:rewrite-uncapped-unit-bgp}}.

No constant growth rates or allocation shares are imposed along these
transitions. The result is local in initial stocks, not in time. In particular,
it concerns equilibrium choices of consumption and research, not convergence
from arbitrary nearby values of all four variables: there are also two
unstable directions. It does not rule out equilibria that leave the local
neighborhood.

\subsection{Growth and distribution without an upper bound}

The uncapped BGP has
\begin{equation}
 \frac{wL}{Y}=(1-\alpha)\omega_L>0,\qquad
 \frac{(r+\delta)K}{Y}=\alpha,\qquad
 \frac{\Pi}{Y}=(1-\alpha)\omega_X-u^*-m^*.
 \label{eq:rewrite-uncapped-unit-distribution}
\end{equation}
Equation~\eqref{eq:rewrite-uncapped-unit-wage-growth} shows that wages
and output per worker grow together at a rate above $\gamma$, even though
labor retains a positive, constant income share. Without an upper bound,
continued growth in AI efficiency sustains this wage-growth premium.
With any upper bound, the unit-elastic equilibria in
Proposition~\ref{prop:rewrite-capped-unit-existence} instead have
$g_{Y/N},g_w\to\gamma$. The equivalence between wage growth above
$\gamma$ and a vanishing labor share across the four regimes in
Proposition~\ref{prop:rewrite-equilibrium-regimes} does not extend to the uncapped economy.

Within the restrictions of Proposition~\ref{prop:rewrite-uncapped-unit-bgp},
increasing the AI weight raises the common growth rate and the interest rate:
\begin{equation}
 \frac{\partial g_{Y/N}^*}{\partial\omega_X}
 =\frac{\partial g_w^*}{\partial\omega_X}
 =\frac{\partial r^*}{\partial\omega_X}
 =\frac{\eta(1-\eta)(n+\gamma)}{(1-\eta-\omega_X)^2}>0.
 \label{eq:rewrite-uncapped-unit-comparative-statics}
\end{equation}
Labor's income share falls across these parameterizations, but remains
constant over time on each BGP. As $\omega_X\to0^+$, the common per-person
growth rate converges to $\gamma$ and $r^*\to\rho+\gamma$. These statements
compare balanced-growth equilibria; they do not describe a transition
from a no-AI economy or establish an uncapped equilibrium for $\sigma>1$.
